\documentclass[../NDCNNAI_main.tex]{subfiles}
\graphicspath{{\subfix{../fig/}}}

\begin{document}
\subsection{N-функции, модуль и норма Люксембурга}

\begin{definition}[N-функция, Young function]
Функция $\theta\colon\mathbb{R}\to[0,+\infty], \theta(0)=0$ такая, 
что она выпуклая, полунепрерывная снизу, чётная, называется N-функцией.
\end{definition}

Например, N-функциями являются:
\begin{itemize}
    \item $\theta_p(s)=\frac{|s|^p}{p} \text{ при }p>1;$ 
    \item $\theta(s)=e^{|s|}-1.$
\end{itemize}

\begin{definition}[Модуль Люксембурга]
Для измеримой функции $f$, заданной на некотором множестве с мерой $(\Omega, \mu)$, 
модуль Люксембурга (ассоциированный с N-функцией $\theta$)
определяется как:
$$\rho_\theta(f)=\int_\Omega\theta(|f(x)|)d\mu(x).$$
\end{definition}

\begin{definition}[Норма Люксембурга]
Под нормой Люксембурга понимается величина:
$$||f||_\theta=\inf\{\lambda>0\colon\rho_\theta\left(\frac{f}{\lambda}\right)\le1\}.$$
\end{definition}

\begin{remark}
Определение корректно, норма Люксембурга действительно является нормой.
\end{remark}

\subsection{Пространства Орлича}
$(\Omega, \mu)$ --- множество с мерой.

\begin{definition}[Пространство Орлича]
Пространством Орлича $L_\theta(\mu)$ называется множество измеримых функций $f$, 
для которых выполняется $||f||_\theta<\infty$.
\end{definition}

\begin{remark}
Некоторые операторы (например, максимальный оператор Харди-Литтлвуда) не ограничены в $L^1$, но ограничены в $L\log^+L$. Это мотивирует выйти за пределы $L^p$, что позволяют делать пространства Орлича с помощью N-функций.
\end{remark}

Примеры пространств Орлича:
\begin{itemize}
    \item Если $\theta(s)=\frac{|s|^p}{p}\text{ при }p>1$, то $L_\theta(\mu)=L^p$.
    \item Если $e^{|s|}-1$, то полученное пространство называется экспоненциальным классом Орлича.
    Возникает в теории вложений Соболева и анализе,
    когда $L^p$-оценок недостаточно (например, соболева функция на $n$-мерной области принадлежит экспоненциальному классу Орлича).
\end{itemize}

\subsection{Обобщённые (Musielak–Orlicz) пространства Орлича}

\begin{definition}[Обобщённая N-функция]
Обобщённая N-функция --- это функция 
$$\Phi\colon\Omega \times\mathbb{R}\to[0,\infty]$$
такая, что:
    \begin{enumerate}
        \item $\Phi(x,\cdot)$ --- N-функция почти всюду для $x\in\Omega$;
        \item Для каждого $t\ge0$ функция $x\mapsto\Phi(x,t)$ --- измерима.
    \end{enumerate}
\end{definition}

\begin{definition}[Обобщённое пространство Орлича]
Обобщённым пространством Орлича называется множество измеримых функций $f$, для которых $\exists\lambda>0$ $\int_\Omega\Phi(x,\frac{f(x)}{\lambda})d\mu(x)<\infty$.
\end{definition}

\begin{remark}
В пространствах Орлича функция $\theta(s)$ имеет однородный рост по всей области.
Обобщённые пространства Орлича позволяют работать с функциями переменного роста,
что позволяет получить, например, переменные пространства Лебега.
\end{remark}

Примеры обобщённых пространств Орлича:
\begin{itemize}
    \item $\Phi(x,t)=\theta(t)\Longrightarrow$ Пространство Орлича;
    \item $\Phi(x,t)=|t|^{p(x)}\Longrightarrow$ Переменное пространство Лебега $(L^{p(\cdot)})$;
    \item $\Phi(x,t)=w(x)|t|^{p}\Longrightarrow$ Взвешенное пространство Лебега $(L^p(w))$.
\end{itemize}

\end{document}